% GENERATED FILE. Edit canonical lessons, then run npm run generate:books.
% canonical-source-hash: fnv1a64:44be703098865f7f
% canonical-lessons: SA-C316-samrat, SA-C316-simhasanam, SA-C316-isvarah, SA-C316-rajyam, SA-C316-sasti

\chapter{Kings and Emperors}
\label{ch:sa-kings-and-emperors}
\hlchaptermodality{316}

\begin{chapteropening}
\textbf{By the end of this chapter:} \emph{I can tell who ruled long ago and how they lived.}

Complete the last lesson of chapter 316: I can tell who ruled long ago and how they lived.
\end{chapteropening}

\section[\texorpdfstring{\sk{सम्राट्} (samrāṭ)}{सम्राट् (samrāṭ)}]{\sk{सम्राट्} (samrāṭ) — an emperor}

\label{lesson:SA-C316-samrat}



\begin{quote}
Before the new one: say the Sanskrit for is alarmed, then the Sanskrit for hurries.
\end{quote}

\subsection*{You'll want to know: \sk{सम्राट्}}
\textbf{\sk{सम्राट्}} — \emph{samrāṭ} — \textquotedblleft{}an emperor\textquotedblright{}.

\begin{quote}
\sk{सम्राट्} \sk{प्रासादे} \sk{वसति।}
\end{quote}

\begin{quote}
\emph{samrāṭ prāsāde vasati.}
\end{quote}

\begin{quote}
The emperor lives in the palace.
\end{quote}

A masculine noun.

\begin{cousinweb}
\textbf{\sk{सम्राट्}} (\emph{samrāṭ}) is built from \emph{sam}, \textquotedblleft{}fully\textquotedblright{}, and the root √rāj, \textquotedblleft{}to shine, rule\textquotedblright{}; MW gives \emph{samrāj} as a universal ruler, a paramount sovereign (MW).
\end{cousinweb}

\begin{grammarlens}[title={What you've built}]
One more word for telling what happened.
\end{grammarlens}

\subsection*{Your turn}
\begin{itemize}
\raggedright
  \item \emph{Say it:} \emph{samrāṭ}
  \item \emph{Say it:} \emph{samrāṭ}, once more
  \item \emph{Say it:} \emph{samrāṭ prāsāde vasati.} — The emperor lives in the palace.
  \item \emph{Recall:} say the Sanskrit for is alarmed, then the Sanskrit for hurries, then say \emph{samrāṭ} again
\end{itemize}

\begin{tcolorbox}[breakable,colback=teal!4,colframe=teal!35!black,title={Before you move on}]
What does \sk{सम्राट्} mean? (\textquotedblleft{}An emperor\textquotedblright{}.) Say it once more.
\end{tcolorbox}
\section[\texorpdfstring{\sk{सिंहासनम्} (siṃhāsanam)}{सिंहासनम् (siṃhāsanam)}]{\sk{सिंहासनम्} (siṃhāsanam) — a throne}

\label{lesson:SA-C316-simhasanam}



\begin{quote}
Before the new one: say the Sanskrit for hurries, then the Sanskrit for an emperor.
\end{quote}

\subsection*{You'll want to know: \sk{सिंहासनम्}}
\textbf{\sk{सिंहासनम्}} — \emph{siṃhāsanam} — \textquotedblleft{}a throne\textquotedblright{}.

\begin{quote}
\sk{सम्राट्} \sk{सिंहासने} \sk{उपविशति।}
\end{quote}

\begin{quote}
\emph{samrāṭ siṃhāsane upaviśati.}
\end{quote}

\begin{quote}
The emperor sits on the throne.
\end{quote}

A neuter noun.

\begin{cousinweb}
\textbf{\sk{सिंहासनम्}} (\emph{siṃhāsanam}) is built from \emph{siṃha}, \textquotedblleft{}a lion\textquotedblright{}, and \emph{āsana}, \textquotedblleft{}a seat\textquotedblright{}; MW gives it as the lion's seat, the king's seat, a throne (MW).
\end{cousinweb}

\begin{grammarlens}[title={What you've built}]
One more word for telling what happened.
\end{grammarlens}

\subsection*{Your turn}
\begin{itemize}
\raggedright
  \item \emph{Say it:} \emph{siṃhāsanam}
  \item \emph{Say it:} \emph{siṃhāsanam}, once more
  \item \emph{Say it:} \emph{samrāṭ siṃhāsane upaviśati.} — The emperor sits on the throne.
  \item \emph{Recall:} say the Sanskrit for hurries, then the Sanskrit for an emperor, then say \emph{siṃhāsanam} again
\end{itemize}

\begin{tcolorbox}[breakable,colback=teal!4,colframe=teal!35!black,title={Before you move on}]
What does \sk{सिंहासनम्} mean? (\textquotedblleft{}A throne\textquotedblright{}.) Say it once more.
\end{tcolorbox}
\section[\texorpdfstring{\sk{ईश्वरः} (īśvaraḥ)}{ईश्वरः (īśvaraḥ)}]{\sk{ईश्वरः} (īśvaraḥ) — a lord; a ruler}

\label{lesson:SA-C316-isvarah}



\begin{quote}
Before the new one: say the Sanskrit for an emperor, then the Sanskrit for a throne.
\end{quote}

\subsection*{You'll want to know: \sk{ईश्वरः}}
\textbf{\sk{ईश्वरः}} — \emph{īśvaraḥ} — \textquotedblleft{}a lord; a ruler\textquotedblright{}.

\begin{quote}
\sk{सम्राट्} \sk{सर्वस्य} \sk{देशस्य} \sk{ईश्वरः} \sk{अस्ति।}
\end{quote}

\begin{quote}
\emph{samrāṭ sarvasya deśasya īśvaraḥ asti.}
\end{quote}

\begin{quote}
The emperor is lord of the whole country.
\end{quote}

A masculine noun.

\begin{cousinweb}
\textbf{\sk{ईश्वरः}} (\emph{īśvaraḥ}) comes from the root √īś, \textquotedblleft{}to own, be master of\textquotedblright{}; MW gives \emph{īśvara} as a master, lord, prince, king, and also God (MW).
\end{cousinweb}

\begin{grammarlens}[title={What you've built}]
One more word for telling what happened.
\end{grammarlens}

\subsection*{Your turn}
\begin{itemize}
\raggedright
  \item \emph{Say it:} \emph{īśvaraḥ}
  \item \emph{Say it:} \emph{īśvaraḥ}, once more
  \item \emph{Say it:} \emph{samrāṭ sarvasya deśasya īśvaraḥ asti.} — The emperor is lord of the whole country.
  \item \emph{Recall:} say the Sanskrit for an emperor, then the Sanskrit for a throne, then say \emph{īśvaraḥ} again
\end{itemize}

\begin{tcolorbox}[breakable,colback=teal!4,colframe=teal!35!black,title={Before you move on}]
What does \sk{ईश्वरः} mean? (\textquotedblleft{}A lord; a ruler\textquotedblright{}.) Say it once more.
\end{tcolorbox}
\section[\texorpdfstring{\sk{राज्यम्} (rājyam)}{राज्यम् (rājyam)}]{\sk{राज्यम्} (rājyam) — a kingdom}

\label{lesson:SA-C316-rajyam}



\begin{quote}
Before the new one: say the Sanskrit for a throne, then the Sanskrit for a lord.
\end{quote}

\subsection*{You'll want to know: \sk{राज्यम्}}
\textbf{\sk{राज्यम्}} — \emph{rājyam} — \textquotedblleft{}a kingdom\textquotedblright{}.

\begin{quote}
\sk{तस्य} \sk{राज्यं} \sk{विशालम्} \sk{अस्ति।}
\end{quote}

\begin{quote}
\emph{tasya rājyaṃ viśālam asti.}
\end{quote}

\begin{quote}
His kingdom is vast.
\end{quote}

A neuter noun.

\begin{cousinweb}
\textbf{\sk{राज्यम्}} (\emph{rājyam}) is built on \emph{rājan}, \textquotedblleft{}a king\textquotedblright{} (\emph{rājā}), with the ending \emph{-ya}; MW gives \emph{rājya} as kingship, sovereignty, and a kingdom (MW).
\end{cousinweb}

\begin{grammarlens}[title={What you've built}]
One more word for telling what happened.
\end{grammarlens}

\subsection*{Your turn}
\begin{itemize}
\raggedright
  \item \emph{Say it:} \emph{rājyam}
  \item \emph{Say it:} \emph{rājyam}, once more
  \item \emph{Say it:} \emph{tasya rājyaṃ viśālam asti.} — His kingdom is vast.
  \item \emph{Recall:} say the Sanskrit for a throne, then the Sanskrit for a lord, then say \emph{rājyam} again
\end{itemize}

\begin{tcolorbox}[breakable,colback=teal!4,colframe=teal!35!black,title={Before you move on}]
What does \sk{राज्यम्} mean? (\textquotedblleft{}A kingdom\textquotedblright{}.) Say it once more.
\end{tcolorbox}
\section[\texorpdfstring{\sk{शास्ति} (śāsti)}{शास्ति (śāsti)}]{\sk{शास्ति} (śāsti) — rules; governs}

\label{lesson:SA-C316-sasti}



\begin{quote}
Before the new one: say the Sanskrit for a lord, then the Sanskrit for a kingdom.
\end{quote}

\subsection*{You'll want to know: \sk{शास्ति}}
\textbf{\sk{शास्ति}} — \emph{śāsti} — \textquotedblleft{}rules; governs\textquotedblright{}.

\begin{quote}
\sk{सम्राट्} \sk{धर्मेण} \sk{राज्यं} \sk{शास्ति।}
\end{quote}

\begin{quote}
\emph{samrāṭ dharmeṇa rājyaṃ śāsti.}
\end{quote}

\begin{quote}
The emperor rules the kingdom justly.
\end{quote}

Past (he …): \sk{सः} \sk{राज्यं} \sk{शासितवान्}, \emph{saḥ rājyaṃ śāsitavān}.

\begin{cousinweb}
\textbf{\sk{शास्ति}} (\emph{śāsti}) comes from the root √śās, \textquotedblleft{}to chastise, correct\textquotedblright{}, which MW also gives as to restrain, rule, govern; MW gives \emph{śāsita} as governed, ruled (MW).
\end{cousinweb}

\begin{grammarlens}[title={What you've built}]
One more word for telling what happened.
\end{grammarlens}

\subsection*{Your turn}
\begin{itemize}
\raggedright
  \item \emph{Say it:} \emph{śāsti}
  \item \emph{Say it:} \emph{śāsti}, once more
  \item \emph{Say it:} \emph{samrāṭ dharmeṇa rājyaṃ śāsti.} — The emperor rules the kingdom justly.
  \item \emph{Recall:} say the Sanskrit for a lord, then the Sanskrit for a kingdom, then say \emph{śāsti} again
\end{itemize}

\begin{tcolorbox}[breakable,colback=teal!4,colframe=teal!35!black,title={Before you move on}]
What does \sk{शास्ति} mean? (\textquotedblleft{}Rules; governs\textquotedblright{}.) Say it once more.
\end{tcolorbox}
